\section*{Glossary}
\label{sec:glossary}
\addcontentsline{toc}{section}{Glossary}

Terms and acronyms used throughout, in order of first appearance. Schedule
names and their parameters are defined in Section~\ref{sec:framework}.

\begin{table}[h]
\centering
\small
\begin{tabular}{p{0.27\textwidth}p{0.66\textwidth}}
\toprule
\multicolumn{2}{l}{\emph{Method and schedule terms}} \\
\midrule
BP & backpropagation: the standard full-depth backward pass. \\
\fw{} (front wheeler) & the proposed schedule: front steps by default,
plateau-triggered bursts of full BP (Algorithm~\ref{alg:fw}). \\
front / trunk / head & the \emph{front} is the last hidden layer (or last
transformer block) plus the output head; the \emph{trunk} is everything
below it; the \emph{head} is the final output projection. \\
front step / full step & a front step back-propagates through the front
only (trunk forwarded without gradients); a full step through everything. \\
burst & a run of consecutive full steps fired by the trigger. \\
trigger, $\tau$ & the plateau detector: a burst fires when the relative
improvement of the training-loss EMA over a check window falls below the
threshold $\tau$. \\
backoff / governor, $\kappa_{\max}$ & if a burst itself improves the loss
by less than $\tau$, the required front-phase length before the next check
doubles (backoff), up to the cap $\kappa_{\max}$ (the governor), which
doubles as the budget dial. \\
deep-step budget & fraction of training steps taken at full depth. \\
backward compute & total backward multiply--accumulates, as a fraction of
\textsc{full}'s; the cost axis used for all cross-schedule comparisons. \\
\midrule
\multicolumn{2}{l}{\emph{Baselines (parameters in
Section~\ref{sec:framework})}} \\
\midrule
\textsc{full} / \textsc{front} & full BP every step / front steps only. \\
\textsc{periodic}-$N$ & one full step every $N$ steps. \\
\lpft{}($s$) & linear probing then fine-tuning: front steps until fraction
$s$ of training, full steps after; one switch. \\
\textsc{ratchet}($r$) & progressive bottom-up freezing: layers freeze
permanently, all-but-front frozen by fraction $r$ of training. \\
\bcd{} & block-coordinate ablation of \fw{}: identical trigger, but bursts
update only the trunk (head frozen). \\
\midrule
\multicolumn{2}{l}{\emph{Optimization and cost terms}} \\
\midrule
EMA & exponential moving average. \\
MAC / GEMM & multiply--accumulate operation / general matrix multiply; we
count backward cost as ${\sim}2$ GEMMs per weight matrix. \\
Adam(W) moments $m, v$ & Adam's per-parameter first and second gradient
moment estimates (EMAs with rates $\beta_1, \beta_2$); AdamW adds decoupled
weight decay. \\
episodic optimizer state & trunk moments that are reset at burst entry (and
warmed up) rather than persisting across inter-burst gaps. \\
dead-ReLU collapse & failure mode where hidden units' pre-activations go
permanently negative, zeroing their gradients and output. \\
Gauss--Southwell rule & greedy block-coordinate selection: update the block
with the largest gradient norm. \\
\midrule
\multicolumn{2}{l}{\emph{Origins-study terms---needed only for Section~\ref{sec:exp1}}} \\
\midrule
supermask & a binary mask over a \emph{frozen, randomly initialized}
network selected so the masked subnetwork performs the task. \\
\ep{} (edge-popup) & gradient-based supermask search: each connection holds
a score; the top-$k$ scores per layer are kept; scores train by BP with a
straight-through estimator. \\
straight-through estimator & treating a non-differentiable operation (here,
top-$k$ selection) as identity during the backward pass. \\
\dfa{} & direct feedback alignment: each layer's error signal is a fixed
random projection of the output error, avoiding weight transport and the
chained backward pass. \\
\rmhebb{} & reward-modulated Hebbian rule: hidden layers receive one global
scalar per sample modulating a local (pre $\times$ post) eligibility trace;
a three-factor rule. \\
eligibility trace & the locally available co-activity term a modulatory
signal multiplies in three-factor learning rules. \\
\bottomrule
\end{tabular}
\end{table}
